\scanpage{28}
\mspage{8}
arithmétiques comme la formulation cohomologique des conjectures de Weil.
Faisons un petit tour d'horizon : à quoi on en est arrivé à ce sujet.

Si $X$ est une variété alg. sur un corps alg. clos $k$,
\deleted{propre sur $k$ pour fixer les idées.}
\inserted{\unsure{on|lui} associe pour tout $\ell$ nb premier}
\uline{$\ell\neq\operatorname{car} k$} des modules sur $\mathbb Z_\ell$
\[
H^i(X,\mathbb Z_\ell)=\varprojlim H^i(X_{\text{ét}},\mathbb Z/\ell^{\unsure{\nu}}\mathbb Z),
\]
dont le comportement pour $\ell$ \deleted{variable} fixé, $X$ variable, est à peu près celui de la cohomologie entière \inserted{ou} rationnelle dans le cadre transcendant.

Ainsi, si $k$ est le corps des complexes, le th. de comparaison nous fournit des isom can, fonctoriels en $X$
\[
\boxed{H^i(X,\mathbb Z_\ell)\simeq H^i(X^{\mathrm{an}},\mathbb Z)\otimes_{\mathbb Z}\mathbb Z_\ell}
\]
Ceci montre que \uline{la connaissance de la cohomologie entière transcendante de $X^{\mathrm{an}}$} (au moins du point de vue additif) \uline{est}
\margin{27}
